\chapter{波导模式与相对论电动力学}

\begin{mapbox}
第一部分仍使用 SI 制，从频域 Maxwell 方程推导 TE、TM、TEM 的纵向本征问题和横向场。第二部分改用 Gaussian 几何单位制并令 $c=1$，以作用量、场张量和 Lorentz 变换呈现相对论结构。两种单位制不在同一推导中混用。
\end{mapbox}

\section{频域 Maxwell 方程与纵横分解（SI）}
均匀、无源、无耗介质中取时间依赖 $e^{-i\omega t}$：
\[
\nabla\times\vb E=i\omega\mu\vb H,\qquad
\nabla\times\vb H=-i\omega\varepsilon\vb E,
\]
\[
\nabla\cdot\vb E=\nabla\cdot\vb H=0.
\]
沿 $z$ 传播的模式写成 $e^{i\beta z}$。分解
\[
\vb E=\vb E_t+\hat{\vb z}E_z,\qquad
\nabla=\nabla_t+i\beta\hat{\vb z}.
\]
定义介质波数 $k=\omega\sqrt{\mu\varepsilon}$ 和
\[
k_c^2=k^2-\beta^2.
\]
把 Maxwell 方程的纵向、横向分量联立消元，可得
\[
\boxed{
\vb E_t=\frac{i}{k_c^2}
\left[\beta\nabla_tE_z
+\omega\mu\,\hat{\vb z}\times\nabla_tH_z\right]},
\]
\[
\boxed{
\vb H_t=\frac{i}{k_c^2}
\left[\beta\nabla_tH_z
-\omega\varepsilon\,\hat{\vb z}\times\nabla_tE_z\right]}.
\]
同时纵向分量都满足二维 Helmholtz 方程：
\[
(\nabla_t^2+k_c^2)E_z=0,\qquad
(\nabla_t^2+k_c^2)H_z=0.
\]
因此只要解出 $E_z$ 或 $H_z$，其余四个横向分量就由上式唯一确定。

\section{TM、TE 与 TEM 的定义}
\[
\begin{array}{c|cc}
&E_z&H_z\\ \hline
\mathrm{TM}&\ne0&0\\
\mathrm{TE}&0&\ne0\\
\mathrm{TEM}&0&0
\end{array}
\]
TM 表示磁场完全横向；TE 表示电场完全横向。波导模式一般并非
$\vb E,\vb H$ 都横向。

理想导体边界上 $\vb E_{\parallel}=0$。对 TM 模，$E_z$ 沿导体壁切向，故
\[
E_z|_{\partial A}=0
\]
是 Dirichlet 条件。对 TE 模，将 $\vb E_t$ 公式投影到壁面切向，可得
\[
\pdv{H_z}{n}\bigg|_{\partial A}=0,
\]
即 Neumann 条件。

\section{矩形波导的 TM 模}
横截面 $0<x<a,\ 0<y<b$。分离变量
$E_z=X(x)Y(y)$：
\[
\frac{X''}{X}+\frac{Y''}{Y}=-k_c^2.
\]
四条边 $E_z=0$，故
\[
E_z=E_0\sin\frac{m\pi x}{a}\sin\frac{n\pi y}{b},
\]
\[
k_c^2=\left(\frac{m\pi}{a}\right)^2+
\left(\frac{n\pi}{b}\right)^2,\qquad m,n=1,2,\ldots.
\]
若任一指标为零，$E_z$ 恒为零，所以不存在 TM$_{m0}$ 或 TM$_{0n}$。
由纵横关系得到
\[
E_x=\frac{i\beta}{k_c^2}\pdv{E_z}{x},\quad
E_y=\frac{i\beta}{k_c^2}\pdv{E_z}{y},
\]
\[
H_x=\frac{i\omega\varepsilon}{k_c^2}\pdv{E_z}{y},\quad
H_y=-\frac{i\omega\varepsilon}{k_c^2}\pdv{E_z}{x}.
\]

\section{矩形波导的 TE 模}
Neumann 条件允许余弦解：
\[
H_z=H_0\cos\frac{m\pi x}{a}\cos\frac{n\pi y}{b},
\]
\[
k_c^2=\left(\frac{m\pi}{a}\right)^2+
\left(\frac{n\pi}{b}\right)^2,
\]
其中 $m,n$ 可以有一个为零，但不能同时为零。横向场为
\[
H_x=\frac{i\beta}{k_c^2}\pdv{H_z}{x},\quad
H_y=\frac{i\beta}{k_c^2}\pdv{H_z}{y},
\]
\[
E_x=-\frac{i\omega\mu}{k_c^2}\pdv{H_z}{y},\quad
E_y=\frac{i\omega\mu}{k_c^2}\pdv{H_z}{x}.
\]
若 $a>b$，最低截止波数是 $k_c=\pi/a$，主模为 TE$_{10}$：
\[
H_z=H_0\cos\frac{\pi x}{a},\qquad
\omega_c=\frac{\pi}{a\sqrt{\mu\varepsilon}}.
\]

\section{截止、色散、速度与波阻抗}
\[
\boxed{\beta^2=k^2-k_c^2
=\mu\varepsilon(\omega^2-\omega_c^2)}.
\]
$\omega>\omega_c$ 时 $\beta$ 为实数，模式传播；低于截止时场沿 $z$ 指数衰减。相速度与群速度
\[
v_p=\frac{\omega}{\beta}
=\frac{v}{\sqrt{1-(\omega_c/\omega)^2}},
\]
\[
v_g=\pdv\omega\beta
=v\sqrt{1-(\omega_c/\omega)^2},
\qquad v=\frac1{\sqrt{\mu\varepsilon}},
\]
所以
\[
\boxed{v_pv_g=v^2}.
\]
$v_p>v$ 不传递信息；包络和能量以 $v_g<v$ 传播。横向场幅度比定义模式阻抗：
\[
Z_{\rm TE}=\frac{\omega\mu}{\beta}
=\frac{Z}{\sqrt{1-(\omega_c/\omega)^2}},
\]
\[
Z_{\rm TM}=\frac{\beta}{\omega\varepsilon}
=Z\sqrt{1-(\omega_c/\omega)^2},
\quad Z=\sqrt{\mu/\varepsilon}.
\]

\section{TEM 模为何需要多连通截面}
若 $E_z=H_z=0$ 且 $k_c=0$，横向场满足二维静电问题
\[
\nabla_t\times\vb E_t=0,\qquad
\nabla_t\cdot\vb E_t=0.
\]
可写 $\vb E_t=-\nabla_t\phi$，其中 $\nabla_t^2\phi=0$。单一空心金属管的整个边界为同一等势面，唯一性定理迫使 $\phi$ 为常数，故不存在非零 TEM 模。同轴线有内外两个独立导体，可以维持电势差，因此支持 TEM。其
\[
\beta=k,\quad v_p=v_g=v,\quad Z_{\rm TEM}=Z
\]
且无截止频率。

\section{模式正交、功率与谐振腔}
二维 Helmholtz 算符在相应边界条件下自伴，不同 $k_c^2$ 的本征函数正交。时间平均功率
\[
P=\frac12\operatorname{Re}\int_A
(\vb E\times\vb H^*)\cdot\hat{\vb z}\,\dd A.
\]
正交性使不同模式的交叉功率积分为零。若在 $z$ 方向再以 PEC 封闭成长度 $d$ 的矩形腔，纵向驻波要求
$k_z=p\pi/d$，谐振频率
\[
\boxed{
\omega_{mnp}=\frac1{\sqrt{\mu\varepsilon}}
\sqrt{\left(\frac{m\pi}{a}\right)^2+
\left(\frac{n\pi}{b}\right)^2+
\left(\frac{p\pi}{d}\right)^2}}.
\]
品质因数 $Q=\omega U/P_{\rm loss}$ 衡量每周期损耗相对储能。

\begin{examplebox}[title=TE$_{10}$ 的截止与色散]
空气波导 $a=2.0\,\mathrm{cm}$。主模截止频率
\[
f_c=\frac{c}{2a}=7.5\,\mathrm{GHz}.
\]
若工作频率 $f=10\,\mathrm{GHz}$，
\[
v_g=c\sqrt{1-(f_c/f)^2}\simeq0.661c,
\quad
v_p=\frac{c^2}{v_g}\simeq1.51c.
\]
乘积仍为 $c^2$，所以超光速相速度没有违反相对论。
\end{examplebox}

\chapter{高斯几何单位制相对论电动力学（光速归一为 1）}

\begin{prereqbox}
本章约定度规 $g_{\mu\nu}=\operatorname{diag}(1,-1,-1,-1)$，
$x^\mu=(t,\vb r)$，并设 $c=1$。时间与长度具有相同量纲，速度是以真空光速为单位的无量纲量，故
$\gamma=(1-v^2)^{-1/2}$。以下公式均为 Gaussian 几何单位制；电荷和场的量纲与 SI 不同。恢复 $c$ 时必须依据量纲和四维定义统一恢复，不能只在个别公式中补写 $c$。
\end{prereqbox}

\section{SI 与 Gaussian 几何制的必要对照}
\[
\begin{array}{c|c|c}
&\mathrm{SI}&\mathrm{Gaussian},\ c=1\\ \hline
\nabla\cdot\vb E&\rho/\varepsilon_0&4\pi\rho\\
\nabla\times\vb B&\mu_0\vb J+c^{-2}\partial_t\vb E&
\ 4\pi\vb J+\partial_t\vb E\\
\nabla\times\vb E&-\partial_t\vb B&-\partial_t\vb B\\
\vb F/q&\vb E+\vb v\times\vb B&
\vb E+\vb v\times\vb B\\
u&(\varepsilon_0E^2+B^2/\mu_0)/2&(E^2+B^2)/(8\pi)
\end{array}
\]
Gaussian 制中 $\vb E,\vb B$ 具有相同量纲；再令 $c=1$ 后，光锥为
$t^2-r^2=0$，四维公式中不再出现人为的 $c$ 因子。表中两列的数值不可直接相等，它们只是同一物理定律在不同单位约定下的形式对照。
本章的三维 Maxwell 方程统一写为
\[
\nabla\cdot\vb E=4\pi\rho,\qquad
\nabla\cdot\vb B=0,
\]
\[
\nabla\times\vb E=-\partial_t\vb B,\qquad
\nabla\times\vb B=4\pi\vb J+\partial_t\vb E.
\]

\section{粒子作用量推出 Lorentz 力}
在 $c=1$ 的 Gaussian 几何单位中，带电粒子的作用量为
\[
\boxed{S_p=-m\int\dd s-q\int A_\mu\dd x^\mu}.
\]
用坐标时间 $t$ 参数化，且取
$A_\mu\dd x^\mu=\phi\,\dd t-\vb A\cdot\dd\vb r$：
\[
L=-m\sqrt{1-v^2}+q\vb A\cdot\vb v-q\phi.
\]
正则动量
\[
\pdv L{\vb v}=\gamma m\vb v+q\vb A.
\]
Euler--Lagrange 方程的第 $i$ 分量：
\[
\dv{}{t}\left(\gamma mv_i+qA_i\right)
=qv_j\partial_iA_j-q\partial_i\phi.
\]
移项并使用
\[
v_j(\partial_iA_j-\partial_jA_i)
=[\vb v\times(\nabla\times\vb A)]_i
\]
以及
$\vb E=-\nabla\phi-\partial_t\vb A$，得到
\[
\boxed{\dv{}t(\gamma m\vb v)
=q\left(\vb E+\vb v\times\vb B\right)}.
\]
$v<1$ 对应有质量粒子低于光速。若恢复普通单位，$\vb v$ 换成
$\vb v/c$，右侧磁力项恢复为 $q\vb v\times\vb B/c$。

\section{场作用量与协变 Maxwell 方程}
定义
\[
A^\mu=(\phi,\vb A),\qquad A_\mu=(\phi,-\vb A),
\]
\[
J^\mu=(\rho,\vb J),\qquad J_\mu=(\rho,-\vb J).
\]
取
\[
S[A]=-\frac1{16\pi}\int F_{\mu\nu}F^{\mu\nu}\dd^4x
-\int J_\mu A^\mu\dd^4x,
\]
\[
F_{\mu\nu}=\partial_\mu A_\nu-\partial_\nu A_\mu.
\]
变分
\[
\delta(F_{\mu\nu}F^{\mu\nu})
=2F^{\mu\nu}(\partial_\mu\delta A_\nu-\partial_\nu\delta A_\mu)
=4F^{\mu\nu}\partial_\mu\delta A_\nu,
\]
其中最后一步使用 $F^{\mu\nu}=-F^{\nu\mu}$。分部积分：
\[
\delta S=
\frac1{4\pi}\int(\partial_\mu F^{\mu\nu})\delta A_\nu\dd^4x
-\int J^\nu\delta A_\nu\dd^4x,
\]
边界项因 $\delta A_\nu|_{\partial\mathcal M}=0$ 消失。任意变分给
\[
\boxed{\partial_\mu F^{\mu\nu}=4\pi J^\nu}.
\]
另一方面 $F=\dd A$ 自动满足
\[
\boxed{\partial_\lambda F_{\mu\nu}
+\partial_\mu F_{\nu\lambda}
+\partial_\nu F_{\lambda\mu}=0}.
\]
二者分别对应非齐次和齐次 Maxwell 方程。

\section{规范不变性与连续性方程}
\[
A_\mu\to A_\mu+\partial_\mu\chi
\]
不改变 $F_{\mu\nu}$。源耦合项变化
\[
\delta S_{\rm int}
=-\int J^\mu\partial_\mu\chi\,\dd^4x
=\int(\partial_\mu J^\mu)\chi\,\dd^4x
\]
（舍去边界项）。要对任意 $\chi$ 保持规范不变，必须
\[
\boxed{\partial_\mu J^\mu=0}
\quad\Longleftrightarrow\quad
\partial_t\rho+\nabla\cdot\vb J=0.
\]
也可直接对场方程取 $\partial_\nu$，因对称导数收缩反对称张量恒为零而得到同一结论。

\section{场张量、Lorentz 变换与不变量}
一种一致约定为
\[
F^{\mu\nu}=
\begin{pmatrix}
0&-E_x&-E_y&-E_z\\
E_x&0&-B_z&B_y\\
E_y&B_z&0&-B_x\\
E_z&-B_y&B_x&0
\end{pmatrix}.
\]
沿 $x$ 方向速度 $\vb v=v\hat{\vb x}$ 的 boost 给
\[
E_x'=E_x,\quad B_x'=B_x,
\]
\[
\vb E_\perp'=\gamma(\vb E_\perp+\vb\beta\times\vb B),
\]
\[
\vb B_\perp'=\gamma(\vb B_\perp-\vb\beta\times\vb E),
\qquad \vb\beta=\vb v,\quad \gamma=(1-v^2)^{-1/2}.
\]
由张量收缩得到两个 Lorentz 不变量：
\[
\boxed{I_1=E^2-B^2},\qquad
\boxed{I_2=\vb E\cdot\vb B}.
\]
若 $I_2=0$ 且 $I_1>0$，存在某参考系使 $\vb B'=0$；若
$I_2=0$ 且 $I_1<0$，存在某系使 $\vb E'=0$。平面真空电磁波满足
$I_1=I_2=0$，不能通过有限 boost 变成纯电场或纯磁场。

\section{能动张量与局域守恒}
Gaussian 几何制电磁能动张量可写为
\[
T^{\mu\nu}=\frac1{4\pi}
\left(-F^{\mu\lambda}F^\nu{}_\lambda
+\frac14g^{\mu\nu}F_{\alpha\beta}F^{\alpha\beta}\right)
\]
（整体号随度规和 $F$ 定义而变，但物理分量固定）。其分量包含
\[
T^{00}=\frac{E^2+B^2}{8\pi},\qquad
T^{0i}=S_i,\qquad
\vb S=\frac1{4\pi}\vb E\times\vb B.
\]
利用 Maxwell 方程和 Bianchi 恒等式可得
\[
\boxed{\partial_\mu T^{\mu\nu}
=-F^{\nu\lambda}J_\lambda}.
\]
右侧是场对物质四维力密度的负值；把物质能动张量加入后，总能量动量局域守恒。

\begin{interviewbox}
\textbf{TE/TM 的回答顺序：}先按纵向分量定义，再说 $E_z,H_z$ 满足横向 Helmholtz 方程；PEC 使 TM 对应 Dirichlet、TE 对应 Neumann；本征值给截止频率，横向场由 Maxwell 方程恢复。TEM 要求 $k_c=0$，单连通空心金属管因唯一性定理不能支持它。

\textbf{相对论电动力学的回答顺序：}四维势生成反对称场张量；场作用量变分给非齐次 Maxwell 方程，$F=\dd A$ 给齐次方程；规范不变性要求电流守恒；$F^{\mu\nu}$ 的 Lorentz 变换统一解释电场和磁场的混合。
\end{interviewbox}

\begin{selftestbox}
\begin{enumerate}
\item 从频域 Maxwell 方程重新推出 $\vb E_t,\vb H_t$。
\item 为什么 TM 模的两个指标都不能为零？
\item 推导 TE 和 TM 模式阻抗。
\item 为什么 TE$_{10}$ 是 $a>b$ 矩形波导的主模？
\item 用唯一性定理说明空心矩形波导没有 TEM。
\item 从粒子作用量逐分量推出 $c=1$ 的 Gaussian 几何制 Lorentz 力。
\item 场作用量中的 $1/16\pi$ 为什么在变分后变成 $1/4\pi$？
\item 用两个场不变量判断能否找到纯电场参考系。
\end{enumerate}
\end{selftestbox}
